% GENERATED FILE. Edit canonical lessons, then run npm run generate:books.
% canonical-source-hash: fnv1a64:fa332af68cb2504c
% canonical-lessons: GU-C100-kaik, GU-C100-kasu, GU-C100-pote, GU-C100-potanu, GU-C100-banne, GU-R100-first-pass-seasons-months-and-days, GU-R100-second-pass-months-this-and-that

\chapter{Something, Anything, Oneself, One's own, Both}
\label{ch:gu-something-anything-oneself-one-s-own-both}
\hlchaptermodality{100}

\begin{chapteropening}
\textbf{By the end of this chapter:} \emph{I can say something, anything, oneself, one's own and both.}

Complete the last lesson of chapter 100: i can say something, anything, oneself, one's own and both.
\end{chapteropening}

\section[kaīk]{\gu{કંઈક} (kaīk) — something}

\label{lesson:GU-C100-kaik}



\begin{quote}
Before the new one: say the Gujarati for every, then the Gujarati for someone.
\end{quote}

\subsection*{You'll want to know: \gu{કંઈક}}
\textbf{\gu{કંઈક}} (\emph{kaīk}) — \textquotedblleft{}something\textquotedblright{}.

\begin{grammarlens}[title={What you've built}]
One more word for this chapter.
\end{grammarlens}

\subsection*{Your turn}
\begin{itemize}
\raggedright
  \item \emph{Say it:} \emph{kaīk}
  \item \emph{Say it:} \emph{kaīk}, once more
  \item \emph{Say it:} say \emph{kaīk}
  \item \emph{Recall:} say the Gujarati for every, then the Gujarati for someone, then say \emph{kaīk} again
\end{itemize}

\begin{tcolorbox}[breakable,colback=teal!4,colframe=teal!35!black,title={Before you move on}]
What does \gu{કંઈક} mean? (\textquotedblleft{}Something\textquotedblright{}.) Say it once more.
\end{tcolorbox}
\section[kaśũ]{\gu{કશું} (kaśũ) — anything}

\label{lesson:GU-C100-kasu}



\begin{quote}
Before the new one: say the Gujarati for someone, then the Gujarati for something.
\end{quote}

\subsection*{You'll want to know: \gu{કશું}}
\textbf{\gu{કશું}} (\emph{kaśũ}) — \textquotedblleft{}anything\textquotedblright{}.

\begin{grammarlens}[title={What you've built}]
One more word for this chapter.
\end{grammarlens}

\subsection*{Your turn}
\begin{itemize}
\raggedright
  \item \emph{Say it:} \emph{kaśũ}
  \item \emph{Say it:} \emph{kaśũ}, once more
  \item \emph{Say it:} say \emph{kaśũ}
  \item \emph{Recall:} say the Gujarati for someone, then the Gujarati for something, then say \emph{kaśũ} again
\end{itemize}

\begin{tcolorbox}[breakable,colback=teal!4,colframe=teal!35!black,title={Before you move on}]
What does \gu{કશું} mean? (\textquotedblleft{}Anything\textquotedblright{}.) Say it once more.
\end{tcolorbox}
\section[pote]{\gu{પોતે} (pote) — oneself}

\label{lesson:GU-C100-pote}



\begin{quote}
Before the new one: say the Gujarati for something, then the Gujarati for anything.
\end{quote}

\subsection*{You'll want to know: \gu{પોતે}}
\textbf{\gu{પોતે}} (\emph{pote}) — \textquotedblleft{}oneself\textquotedblright{}.

\begin{grammarlens}[title={What you've built}]
One more word for this chapter.
\end{grammarlens}

\subsection*{Your turn}
\begin{itemize}
\raggedright
  \item \emph{Say it:} \emph{pote}
  \item \emph{Say it:} \emph{pote}, once more
  \item \emph{Say it:} say \emph{pote}
  \item \emph{Recall:} say the Gujarati for something, then the Gujarati for anything, then say \emph{pote} again
\end{itemize}

\begin{tcolorbox}[breakable,colback=teal!4,colframe=teal!35!black,title={Before you move on}]
What does \gu{પોતે} mean? (\textquotedblleft{}Oneself\textquotedblright{}.) Say it once more.
\end{tcolorbox}
\section[potānũ]{\gu{પોતાનું} (potānũ) — one's own}

\label{lesson:GU-C100-potanu}



\begin{quote}
Before the new one: say the Gujarati for anything, then the Gujarati for oneself.

Say \textquotedblleft{}I speak Gujarati.\textquotedblright{} (\emph{Hũ gujarātī bolũ chhũ}.)
\end{quote}

\subsection*{You'll want to know: \gu{પોતાનું}}
\textbf{\gu{પોતાનું}} (\emph{potānũ}) — \textquotedblleft{}one's own\textquotedblright{}.

\begin{grammarlens}[title={What you've built}]
One more word for this chapter.
\end{grammarlens}

\subsection*{Your turn}
\begin{itemize}
\raggedright
  \item \emph{Say it:} \emph{potānũ}
  \item \emph{Say it:} \emph{potānũ}, once more
  \item \emph{Say it:} say \emph{potānũ}
  \item \emph{Recall:} say the Gujarati for anything, then the Gujarati for oneself, then say \emph{potānũ} again
\end{itemize}

\begin{tcolorbox}[breakable,colback=teal!4,colframe=teal!35!black,title={Before you move on}]
What does \gu{પોતાનું} mean? (\textquotedblleft{}One's own\textquotedblright{}.) Say it once more.
\end{tcolorbox}
\section[banne]{\gu{બંને} (banne) — both}

\label{lesson:GU-C100-banne}



\begin{quote}
Before the new one: say the Gujarati for oneself, then the Gujarati for one's own.
\end{quote}

\subsection*{You'll want to know: \gu{બંને}}
\textbf{\gu{બંને}} (\emph{banne}) — \textquotedblleft{}both\textquotedblright{}.

\begin{grammarlens}[title={What you've built}]
One more word for this chapter.
\end{grammarlens}

\subsection*{Your turn}
\begin{itemize}
\raggedright
  \item \emph{Say it:} \emph{banne}
  \item \emph{Say it:} \emph{banne}, once more
  \item \emph{Say it:} say \emph{banne}
  \item \emph{Recall:} say the Gujarati for oneself, then the Gujarati for one's own, then say \emph{banne} again
\end{itemize}

\begin{tcolorbox}[breakable,colback=teal!4,colframe=teal!35!black,title={Before you move on}]
What does \gu{બંને} mean? (\textquotedblleft{}Both\textquotedblright{}.) Say it once more.
\end{tcolorbox}
\section[(dialogue)]{Seasons, months and days — a first pass}

\label{lesson:GU-R100-first-pass-seasons-months-and-days}



\begin{quote}
Say the words for one's own and both.
\end{quote}

\subsection*{Your turn}
Say these aloud:

\begin{itemize}
\raggedright
  \item midnight, dawn, the day after tomorrow, a year, a moment
  \item a season, summer, winter, the monsoon, spring
  \item autumn, early, every day, a date, right now
  \item still, immediately, January, February, March
\end{itemize}

\begin{tcolorbox}[breakable,colback=teal!4,colframe=teal!35!black,title={Before you move on}]
Midnight? (\textbf{\gu{મધરાત}}.) Dawn? (\textbf{\gu{પરોઢ}}.) The day after tomorrow? (\textbf{\gu{પરમદિવસ}}.) A year? (\textbf{\gu{વરસ}}.) A moment? (\textbf{\gu{પળ}}.) A season? (\textbf{\gu{મોસમ}}.) Summer? (\textbf{\gu{ઉનાળો}}.) Winter? (\textbf{\gu{શિયાળો}}.) The monsoon? (\textbf{\gu{ચોમાસું}}.) Spring? (\textbf{\gu{વસંત}}.) Autumn? (\textbf{\gu{પાનખર}}.) Early? (\textbf{\gu{વહેલું}}.) Every day? (\textbf{\gu{રોજ}}.) A date? (\textbf{\gu{તારીખ}}.) Right now? (\textbf{\gu{હમણાં}}.) Still? (\textbf{\gu{હજી}}.) Immediately? (\textbf{\gu{તરત}}.) January? (\textbf{\gu{જાન્યુઆરી}}.) February? (\textbf{\gu{ફેબ્રુઆરી}}.) March? (\textbf{\gu{માર્ચ}}.)
\end{tcolorbox}
\section[(dialogue)]{Months, this and that — a second pass}

\label{lesson:GU-R100-second-pass-months-this-and-that}



\begin{quote}
Say the words for April and May.
\end{quote}

\subsection*{Your turn}
Say these aloud:

\begin{itemize}
\raggedright
  \item April, May, June, July, September
  \item November, Kartak, Magshar, Phagan, Chaitra
  \item Vaishakh, Jeth, Shravan, Bhadarvo, Aso
  \item that, like this, like that, every, someone
  \item something, anything, oneself, one's own, both
\end{itemize}

\begin{tcolorbox}[breakable,colback=teal!4,colframe=teal!35!black,title={Before you move on}]
April? (\textbf{\gu{એપ્રિલ}}.) May? (\textbf{\gu{મે}}.) June? (\textbf{\gu{જૂન}}.) July? (\textbf{\gu{જુલાઈ}}.) September? (\textbf{\gu{સપ્ટેમ્બર}}.) November? (\textbf{\gu{નવેમ્બર}}.) Kartak? (\textbf{\gu{કારતક}}.) Magshar? (\textbf{\gu{માગશર}}.) Phagan? (\textbf{\gu{ફાગણ}}.) Chaitra? (\textbf{\gu{ચૈત્ર}}.) Vaishakh? (\textbf{\gu{વૈશાખ}}.) Jeth? (\textbf{\gu{જેઠ}}.) Shravan? (\textbf{\gu{શ્રાવણ}}.) Bhadarvo? (\textbf{\gu{ભાદરવો}}.) Aso? (\textbf{\gu{આસો}}.) That? (\textbf{\gu{પેલું}}.) Like this? (\textbf{\gu{આવું}}.) Like that? (\textbf{\gu{એવું}}.) Every? (\textbf{\gu{દરેક}}.) Someone? (\textbf{\gu{કોઈ}}.) Something? (\textbf{\gu{કંઈક}}.) Anything? (\textbf{\gu{કશું}}.) Oneself? (\textbf{\gu{પોતે}}.) One's own? (\textbf{\gu{પોતાનું}}.) Both? (\textbf{\gu{બંને}}.)
\end{tcolorbox}
